% ============================== UNIT 1 =====================================
\unitheading{1}{Classical description of the atom}

\textit{Handout chapter 1, sections 1.1--1.3.} The aim of this unit is to know what an atom is made of,
how we found out, and how to describe a nucleus with the two numbers $Z$ and $A$.

\partheading{Concepts in detail}

\sect{1.1}{Imaging atoms: from an idea to an object}
\begin{concept}[Atom.]
The smallest particle of a chemical element that keeps the identity of that element. The word comes
from the Greek \emph{atomos}, ``that cannot be cut''.
\end{concept}
\begin{itemize}
  \item \textbf{Democritus} (5th century BC): matter cannot be divided forever; there must be
        indivisible building blocks, the atoms. This was a philosophical idea with no experiment behind it.
  \item \textbf{Dalton} (1808): the first \emph{scientific} atomic theory. Each element is made of
        one type of atom with its own mass, and a given compound always contains the same ratio of
        atoms (for example water always has 2 H for 1 O). This explains the law of definite proportions.
  \item \textbf{Scanning tunnelling microscopy} (STM, 1980s): a sharp metal tip is moved a few
        tenths of a nanometre above a surface; the tiny current that ``tunnels'' between tip and surface
        depends very strongly on the distance, so the surface can be mapped atom by atom. The same tip
        can push atoms around: Figure~1-1 of the handout is an ellipse of cobalt atoms placed one by one
        on a copper surface. The atom went from a hypothesis to a real object that can be seen and moved.
\end{itemize}

\sect{1.2}{Internal structure: the atom is not indivisible}
\subsect{1.2.1\quad The electron (Thomson, 1897)}
A high voltage between two electrodes in a tube containing a gas at low pressure produces a beam
coming out of the negative electrode (the cathode): a \emph{cathode ray}.
\begin{itemize}
  \item The beam is deflected by an electric field \emph{towards the positive plate}: it carries a
        \textbf{negative} charge.
  \item The deflection is the same whatever the gas and the metal of the electrodes: the particles are
        a component of \emph{all} atoms.
  \item From the deflection in electric and magnetic fields, Thomson measured the charge-to-mass ratio
        $e/m_e \approx \SI{1.76e11}{\coulomb\per\kilo\gram}$, about 1800 times larger than for the lightest
        ion (H$^+$). The particle is therefore much lighter than an atom: the \textbf{electron}.
\end{itemize}
Consequence: atoms contain electrons, and since matter is neutral they must also contain positive
charge. Thomson imagined the ``plum pudding'' model: electrons embedded in a uniform positive sphere.

\subsect{1.2.2\quad The nucleus (Rutherford, 1911)}
$\alpha$-particles (He$^{2+}$ ions, positive and heavy) are fired at a very thin gold foil, and a
detector all around the foil records where they go.
\begin{center}
\small
\begin{tabular}{@{}>{\raggedright\arraybackslash}p{62mm} >{\raggedright\arraybackslash}p{88mm}@{}}
\toprule
\textbf{Observation} & \textbf{Conclusion}\\ \midrule
Almost all $\alpha$-particles pass straight through. & The atom is mostly \textbf{empty space}.\\
A few are deflected by large angles. & They pass close to a small region of concentrated
\textbf{positive} charge that repels them.\\
Very rarely (about 1 in 10\,000) one bounces back. & That region is tiny but contains almost
\textbf{all the mass} of the atom: the \textbf{nucleus}.\\ \bottomrule
\end{tabular}
\end{center}
The plum-pudding model cannot explain the bounce-backs: a diffuse positive charge would never stop a
fast, heavy $\alpha$-particle. Rutherford's model: a tiny, dense, positive nucleus surrounded by
electrons at comparatively huge distances.

\subsect{1.2.3\quad Sizes, charges and masses}
\begin{itemize}
  \item Size of an atom: about \SI{100}{\pico\metre} $= \SI{e-10}{\metre}$ ($1\ \mathrm{pm} =
        \SI{e-12}{\metre}$; the older unit $1$~\AA{} $= \SI{100}{\pico\metre}$).
  \item Size of a nucleus: $10^{-3}$ to $10^{-2}$~pm, that is 1 to 10~femtometres
        ($1\ \mathrm{fm} = \SI{e-15}{\metre}$). The atom is $10^4$ to $10^5$ times larger than its nucleus:
        if the nucleus were a marble of 1~cm, the atom would be a sphere of about 1~km.
  \item \textbf{Millikan} (1909, oil-drop experiment): every charge he measured was a whole multiple of
        the same value, the \textbf{elementary charge} $e = \SI{1.6e-19}{\coulomb}$. Electron: $-e$;
        proton: $+e$; neutron: 0.
  \item Masses (Table 1-1): $m_e = \SI{9.1094e-31}{\kilo\gram}$, $m_p = \SI{1.6727e-27}{\kilo\gram}$,
        $m_n = \SI{1.6749e-27}{\kilo\gram}$. The proton and neutron have almost the same mass, and are
        about \textbf{1836 times} heavier than the electron. Almost all the mass of an atom is in its nucleus.
\end{itemize}
\begin{concept}[Atomic mass unit (amu, also written u or Da).]
Exactly $\frac1{12}$ of the mass of one unbound carbon-12 atom, at rest and in its ground state:
$1\ \mathrm{amu} = \SI{1.6605e-27}{\kilo\gram}$. It is chosen so that protons and neutrons have
masses close to 1~amu, and it links to the mole: $1\ \mathrm{amu}\times N_A = \SI{1}{\gram\per\mole}$.
\end{concept}
\begin{worked}[: converting Table 1-1 into amu]
$m_p/(1\ \mathrm{amu}) = 1.6727/1.6605 = 1.0073$;\quad $m_e/(1\ \mathrm{amu}) = 9.1094\times10^{-31}/1.6605\times10^{-27}
= 5.49\times10^{-4} = 1/1823$;\quad gold atom: $3.2707\times10^{-25}/1.6605\times10^{-27} = 196.97$.
(For the hydrogen atom the same division gives 1.0078, which rounds to 1.008.)
\end{worked}

\sect{1.3}{Chemical element and isotopes}
\begin{concept}[Atomic number $Z$ and mass number $A$.]
$Z$ = number of protons in the nucleus (also the number of electrons in the neutral atom).
$A$ = number of nucleons (protons + neutrons). The number of neutrons is $N = A - Z$.
A nucleus (or atom) is written $\nuc{A}{Z}{X}$, where X is the symbol of the element, for example
$\nuc{12}{6}{C}$, $\nuc{35}{17}{Cl}$, $\nuc{56}{26}{Fe}$.
\end{concept}
\begin{concept}[Chemical element.]
The set of all atoms and ions that have the same atomic number $Z$. $Z$ alone fixes the element:
every atom with 6 protons is carbon, whatever its number of neutrons or electrons.
\end{concept}
\begin{concept}[Isotopes.]
Atoms of the same element (same $Z$) with different numbers of neutrons, hence different $A$.
Example: $\nuc{12}{6}{C}$ (6 neutrons), $\nuc{13}{6}{C}$ (7) and $\nuc{14}{6}{C}$ (8). Isotopes have the
same electrons, hence the same chemistry; only their mass (and nuclear stability) differ.
\end{concept}
\begin{concept}[Ions.]
An atom that has lost or gained electrons. The nucleus is unchanged. A cation $\mathrm{X}^{q+}$ has
$Z-q$ electrons; an anion $\mathrm{X}^{q-}$ has $Z+q$ electrons. Species with the same number of electrons
are called \textbf{isoelectronic} (for example $\ce{F-}$, Ne, $\ce{Na+}$, $\ce{Mg^2+}$ all have 10 electrons).
\end{concept}
\begin{trap}
$A$ is a \emph{count} of nucleons (an integer), not the atomic mass. The mass of $\nuc{35}{17}{Cl}$ is
34.97~amu, and the atomic mass of chlorine found in tables (35.45) is an average over its isotopes
(Unit 5). Also, changing the number of electrons makes an ion, not an isotope.
\end{trap}

% ------------------------------------------------------------------ recap
\clearpage
\begin{recap}{Unit 1: classical description of the atom}
\recapsub{Timeline}
\begin{tabular}{@{}l l l@{}}
5th c. BC & Democritus & indivisible atoms (philosophy)\\
1808 & Dalton & each element has its own atoms and mass; fixed ratios\\
1897 & Thomson & electron: cathode rays, negative, $e/m_e$ very large\\
1909 & Millikan & charge is quantised: $e = \SI{1.6e-19}{\coulomb}$\\
1911 & Rutherford & gold foil: tiny, dense, positive nucleus; the atom is mostly empty\\
1980s & STM & individual atoms imaged and moved
\end{tabular}

\recapsub{Particles}
\begin{tabular}{@{}l c c c@{}}
\toprule
 & charge & mass (kg) & mass (amu)\\ \midrule
electron & $-e$ & $9.1094\times10^{-31}$ & $1/1823 \approx 0.00055$\\
proton & $+e$ & $1.6727\times10^{-27}$ & 1.0073\\
neutron & 0 & $1.6749\times10^{-27}$ & 1.0087\\ \bottomrule
\end{tabular}\hfill
\begin{minipage}[t]{40mm}\small
$1\ \mathrm{amu} = \SI{1.6605e-27}{\kilo\gram}$\\ $= \frac1{12}\,m(\iso{12}{C})$\\[2pt]
$m_p/m_e \approx 1836$\\[2pt]
$1\ \mathrm{amu}\times N_A = 1\ \mathrm{g\,mol^{-1}}$
\end{minipage}

\recapsub{Orders of magnitude}
Atom $\sim\SI{100}{\pico\metre} = \SI{e-10}{\metre}$; nucleus $\sim 10^{-3}$--$10^{-2}$~pm (1--10~fm);
ratio $10^4$--$10^5$. The nucleus holds $>99.9\,\%$ of the mass.

\recapsub{Notation and counting}
$\nuc{A}{Z}{X}$:\quad protons $=Z$,\quad neutrons $N = A-Z$,\quad electrons $= Z - q$ for charge $q$
(so $Z+|q|$ for an anion).\\
Element $\Leftrightarrow$ same $Z$.\quad Isotopes $\Leftrightarrow$ same $Z$, different $A$.\quad
Isoelectronic $\Leftrightarrow$ same number of electrons.

\recapsub{You should be able to}
\begin{itemize}
  \item describe the Thomson and Rutherford experiments and say what each observation proves;
  \item convert between pm, nm, \AA{}, fm and m, and between kg and amu;
  \item give $Z$, $A$, $N$ and the electron count of any atom or ion written $\nuc{A}{Z}{X}^{q\pm}$;
  \item recognise isotopes and isoelectronic species.
\end{itemize}

\recapsub{Watch out}
\begin{itemize}
  \item Cathode rays bend \emph{towards} the positive plate (they are negative).
  \item A cation has \emph{fewer} electrons than protons; an anion has \emph{more}.
  \item $A$ is an integer count; masses in amu are not integers (except $\iso{12}{C}$, by definition).
\end{itemize}
\end{recap}
